\documentclass[11pt,a4paper]{article}

% --- Packages ---
\usepackage[utf8]{inputenc}
\usepackage[T1]{fontenc}
\usepackage[english]{babel}
\usepackage{amsmath, amssymb, amsthm}
\usepackage{graphicx}
\usepackage{geometry}
\usepackage{xcolor}
\usepackage{hyperref}
\usepackage{algorithm}
\usepackage{algpseudocode}
\usepackage{subcaption}
\usepackage{listings}

\geometry{margin=2.5cm}

\newtheorem{theorem}{Theorem}
\newtheorem{definition}{Definition}
\newtheorem{remark}{Remark}
\newtheorem{example}{Example}

% --- Custom Code Style ---
\lstset{
    basicstyle=\ttfamily\small,
    breaklines=true,
    frame=single,
    backgroundcolor=\color{gray!5},
    keywordstyle=\color{blue},
    commentstyle=\color{green!60!black},
    stringstyle=\color{purple}
}

% --- Macros ---
\newcommand{\R}{\mathbb{R}}

\newcommand{\commentforimage}[1]{\text{\color{red} [Image Placeholder]: #1}}



\title{A Domain-Specific Language and Web Interface for Geometric Simplicial Complexes in PolyLogicA}
\author{Simone Riccio}
\date{\today}

\begin{document}

\maketitle

\begin{abstract}
    This report presents the design and implementation of a novel Domain-Specific Language (DSL) and web interface tailored for the construction and manipulation of geometric simplicial complexes. Developed to interface directly with the PolyLogicA branch of the VoxLogicA project, this tool allows users to define topological spaces in Euclidean coordinates, apply geometric transformations, and evaluate spatial logic queries. We detail the language grammar, the abstract syntax tree evaluation, and the topological serialization process required to generate the Hasse diagram (poset) representations necessary for spatial model checking.
\end{abstract}

\section{Introduction}

\section{Theoretical Background}

\subsection{Geometry and Simplicial Complexes}

\begin{definition} [Simplex]
    Given a set of $k+1$ affinely independent points in $\R$ denoted as $\left\{v_0, v_1, \ldots, v_k \right\}$, the $k$-simplex $\sigma$ spanned by these points is defined as the convex hull:
    \[
    \sigma = \left[v_0, v_1, \ldots, v_k\right] = \left\{ \sum_{i=0}^{k} \lambda_i v_i \mid \lambda_i \geq 0, \sum_{i=0}^{k} \lambda_i = 1 \right\}.
    \]
\end{definition}

\begin{definition} [Simplicial Complex]
    A simplicial complex $K$ is a finite collection of simplices such that:
    \begin{enumerate}
        \item Every face of a simplex in $K$ is also in $K$.
        \item The intersection of any two simplices in $K$ is either empty or a face of both simplices.
    \end{enumerate}
\end{definition}

\commentforimage{Immagini di simplessi e complessi simpliciali di varie dimensioni}

\begin{remark}
    Simplicial complexes provide a combinatorial representation of topological spaces, allowing for the application of algebraic topology techniques and spatial logic reasoning.
\end{remark}

\begin{definition} [Abstract Simplicial Complex]
    An abstract simplicial complex is a set $K$ of finite subsets of a vertex set $V$ such that if $\sigma \in K$ and $\tau \subseteq \sigma$, then $\tau \in K$. The elements of $K$ are called simplices, and the elements of $V$ are called vertices.
\end{definition}

\begin{definition} [Dimension]
    The dimension of a simplex $\sigma = \left[v_0, v_1, \ldots, v_k\right]$ is defined as $k$. The dimension of a simplicial complex $K$ is the maximum dimension of its simplices:
    \[
    \dim(K) = \max\{\dim(\sigma) \mid \sigma \in K\}.
    \]
\end{definition}

\begin{definition} [Downward Closure]
    The downward closure of a simplex $\sigma$ is the set of all its faces:
    \[
    \text{C}(\sigma) = \{\tau \subseteq \sigma \mid \tau \text{ is a face of } \sigma\}.
    \]
\end{definition}

\subsection{Posets and Hasse Diagrams}

\begin{definition} [Partially Ordered Set (Poset)]
    A partially ordered set (poset) is a set $P$ equipped with a binary relation $\leq$ that is reflexive, antisymmetric, and transitive.
\end{definition}

\begin{definition} [Hasse Diagram]
    A Hasse diagram is a graphical representation of a finite poset, where each element is represented as a vertex, and edges are drawn to indicate the covering relations. An edge from vertex $a$ to vertex $b$ is drawn if $a < b$ and there is no $c$ such that $a < c < b$. The diagram is typically drawn so that if $a < b$, then $b$ appears higher than $a$ in the diagram.
\end{definition}

\begin{remark}
    The set of simplices in a simplicial complex, ordered by inclusion, forms a poset. The Hasse diagram of this poset visually represents the inclusion relationships among the simplices, which is crucial for spatial model checking in PolyLogicA.
\end{remark}

\commentforimage{Immagini di un complesso simpliciale a sinistra e il suo diagramma di Hasse a destra}

\begin{remark} 
    The transition from geometric simplicial complexes to partially ordered sets (posets) is the core computational step for interfacing with spatial model checkers like PolyLogicA. While a simplex is defined by its coordinates in Euclidean space, its logical existence in the model checker is defined by its position in a hierarchy of inclusion. 
\end{remark}

\subsection{Spatial Logic}

Broadly defined, a spatial logic is any formal language interpreted over a class of structures featuring goemetrical entities and relations. \cite{spatial-logic} provides a comprehensive overview of the various spatial logics, their semantics, and applications.

The origin of this field lies in the topological approach to modal logic initiated by Alfred Tarski, who recognized that modal operators could represent topological notions such as interior and closure.

Modern developments in spatial logic extend these ideas to discrete spatial structures (like graphs or images) using a generalized topological space called a \textbf{closure space} \cite{voxlogica} 
\begin{definition} [Closure Space]
    A closure space is a pair $(X, \mathcal{C})$ where $X$ is a set and $\mathcal{C}: \mathcal{P}(X) \to \mathcal{P}(X)$ is a closure operator satisfying the following properties for all $A, B \subseteq X$:
    \begin{enumerate}
        \item $\mathcal{C}(\emptyset) = \emptyset$
        \item $A \subseteq \mathcal{C}(A)$ (Extensivity)
        \item If $A, B \subseteq X$ and $\mathcal{C}(A \cup B) = \mathcal{C}(A) \cup \mathcal{C}(B)$ 
    \end{enumerate}
\end{definition}

\begin{remark}
    In any topological space $X$ we can define $\mathcal{C}(A)$ as the topological closure of $A$, which is the smallest closed set containing $A$. This closure operator satisfies the properties of a closure operator, making $(X, \mathcal{C})$ a closure space. 

    The closure operator in this case has one additional property: it is idempotent, meaning that $\mathcal{C}(\mathcal{C}(A)) = \mathcal{C}(A)$ for all $A \subseteq X$. This property is not required for a closure space, but it is a defining characteristic of topological spaces.
\end{remark}

\begin{example} [A closure space which is not a topological space]
    Consider for a oriented graph $G = (V, E)$ the closure operator $\mathcal{C}$ defined as follows: for any $A \subseteq V$, $\mathcal{C}(A)$ is the set of all vertices in $V$ that can be reached from any vertex in $A$ by a directed path in $G$. This closure operator satisfies the properties of a closure operator, making $(V, \mathcal{C})$ a closure space.

    However, this closure operator is not necessarily idempotent. For example, if $A$ consists of a single vertex with outgoing edges to other vertices, then $\mathcal{C}(A)$ will include those reachable vertices. Applying $\mathcal{C}$ again may yield additional vertices that are reachable from the newly included vertices, thus $\mathcal{C}(\mathcal{C}(A))$ may be strictly larger than $\mathcal{C}(A)$. Therefore, $(V, \mathcal{C})$ is a closure space but not a topological space.

    If the graph is undirected, then the closure operator becomes idempotent, and $(V, \mathcal{C})$ is a topological space. 
\end{example}

Ciancia et al. define a spatial logic for closure spaces for reasoning about spatial properties of closure spaces.
\begin{definition} [Spatial Logic for Closure Spaces] \cite{voxlogica}
    In a closure space $(X, \mathcal{C})$, the syntax of the spatial logic is defined as follows:
    \[
    \phi ::= p \mid \neg \phi \mid \phi_1 \wedge \phi_2 \mid \mathcal{N} \phi \mid \phi_1 \mathcal{U} \phi_2
    \]
    where $p$ is a propositional variable, $\neg$ is negation, $\wedge$ is conjunction, $\mathcal{N}$ is the "next" operator, and $\mathcal{U}$ is the "until" operator. The semantics of the logic is defined in terms of the closure operator $\mathcal{C}$ and the satisfaction relation $\models$ as follows:
    \begin{itemize}
        \item $(X, \mathcal{C}), x \models p$ if $x \in V(p)$, where $V(p)$ is the set of points in $X$ where $p$ is true.
        \item $(X, \mathcal{C}), x \models \neg \phi$ if $(X, \mathcal{C}), x \not\models \phi$.
        \item $(X, \mathcal{C}), x \models \phi_1 \wedge \phi_2$ if $(X, \mathcal{C}), x \models \phi_1$ and $(X, \mathcal{C}), x \models \phi_2$.
        \item $(X, \mathcal{C}), x \models \mathcal{N} \phi$ if there exists $y \in \mathcal{C}(\{x\})$ such that $(X, \mathcal{C}), y \models \phi$.
        \item $(X, \mathcal{C}), x \models \phi_1 \mathcal{U} \phi_2$ if there exists a finite sequence of points $x_0, x_1, \ldots, x_n \in X$ such that $x_0 = x$, $(X, \mathcal{C}), x_n \models \phi_2$, and for all $0 \leq i < n$, $(X, \mathcal{C}), x_i \models \phi_1$.
    \end{itemize}
\end{definition}

We will not go into the details of the semantics of the logic, but we will focus on how to represent geometric simplicial complexes in a way that allows us to reason about their spatial properties using this logic.


% \section{Geometric DSL: Syntax and Semantics}

% Prova

% \section{Implementation Details}

% Prova

% \section{Integration with PolyLogicA}

% Prova

% \section{Examples and Model Checking}


% \section{Conclusion and Future Work}

\begin{thebibliography}{9}
    \bibitem{spatial-logic}
    Marco Aiello, Ian Pratt-Hartmann, and Johan van Benthem,
    \textit{Handbook of Spatial Logic},
    Springer, 2007.
    
    \bibitem{voxlogica}
    Vincenzo Ciancia et al.,
    \textit{VoxLogicA: A Spatial Model Checker}.
\end{thebibliography}
    
\end{document}